\documentclass[10pt,a4paper]{amsart}
\usepackage[margin=19mm]{geometry}
\usepackage{amsmath,amssymb,mathtools}
\usepackage[scr=rsfs,cal=euler]{mathalfa}
\usepackage{xfrac}
\usepackage[unicode,hidelinks,bookmarks=false]{hyperref}
\newcommand{\ie}{i.e.\ }
\newcommand{\cf}{cf.\ }
\newcommand{\resp}{respectively\ }
\newcommand{\Subsection}{\subsection}
\providecommand{\nobreakdash}{\nobreak\hbox{-}}
\setlength{\emergencystretch}{2em}
\multlinegap0pt
\usepackage{amssymb}
\usepackage{enumerate}
\usepackage{dsfont}
\usepackage{bm}
\usepackage{mathtools}
\usepackage{algorithm}\usepackage{algpseudocode}
\usepackage[shortlabels]{enumitem}
\newtheorem{assumption}{Assumption}
\newtheorem{lemma}{Lemma}
\def \P{\mathbb{P}}
\def \R{\mathbb{R}}
\def \dd{{\rm d}}
\title{Exploratory optimal stopping: A singular control formulation}
\author{Jodi Dianetti, Giorgio Ferrari, Renyuan Xu}
\date{Source: arXiv:2408.09335v3 (CC BY 4.0)}
\begin{document}
\maketitle

\noindent\emph{Source and licence.} Exploratory optimal stopping: A singular control formulation. Version arXiv:2408.09335v3 (CC BY 4.0), distributed by arXiv under the Creative Commons Attribution 4.0 International licence, http://creativecommons.org/licenses/by/4.0/, as recorded in the arXiv licence metadata for this exact version and shown on the abstract page https://arxiv.org/abs/2408.09335v3. Full licence notice: \texttt{licenses/arxiv-2408.09335-CC-BY-4.0.txt}.\par\smallskip

\noindent\textit{Editorial scope.} Counted unit: the article's lemma times (source0.tex lines 1022--1025). The article's complete original proof of it is delivered with it (source0.tex lines 1026--1056, one \texttt{proof} environment of 31 lines). No mathematics is written, repaired or re-derived: the statement and the proof are the article's own lines, in the article's own notation, and no retained line is substituted or re-flowed.  Retained uncounted as the source-local context the proof needs: lines 853--906; lines 996--1004; lines 1133--1139; lines 1146--1153.  Nothing was omitted from any retained span, so the package carries no omission record.  No retained line contains a citation, so the wrapper carries no bibliography and no bibliography entry.  No retained line contains a figure environment, an image-inclusion command or drawing code, so no image is delivered and the package declares no asset dependency.  Editorial formatting only, in addition: an AMS \texttt{amsart} preamble that restates the article's own declarations and macros used by the retained lines, namely the environments \texttt{assumption}, \texttt{lemma} on separate counters, together with the standard packages those lines need.  The retained environments print in this document as Assumption 1, Lemma 1 in order of appearance, whereas the article prints the same items with the numbering of its own full document.  The complete native source of the article as distributed in the arXiv e-print for this exact version is bundled unchanged as \texttt{sources/163909/source0.tex}.
\begin{assumption}
    \label{assumption general}
    The following conditions hold true: 
    \begin{enumerate}
    \item\label{condition b sigma} There exists a constant $L>0$ such that, for $\phi = b, \sigma$, we have
    $$ \begin{aligned}
         | \phi (x)| & \leq L (1+|x|),\\ 
        |\phi (\bar x) - \phi (x)|  & \leq L |\bar x - x|, \\
        | \phi ( \delta \bar x + (1-\delta)x) - \delta \phi(\bar x) - (1-\delta) \phi ( x) |  &\leq  L \delta (1-\delta)  |\bar{x}-x|^2,
    \end{aligned}
    $$
    for any $x, \bar x \in \mathbb R ^n$, $\delta \in [0,1]$. 
    \item\label{condition costs} There exist $p\geq 2$ and $K>0$ such that, for $\phi = G, \pi$, we have
    $$
    \begin{aligned}
    |\phi(x)| &\leq K (1+|x|^p), \\
    |\phi (\bar x) - \phi (x) | &\leq K (1+|\bar x|^{p-1} + |x|^{p-1})|\bar x -x |, \\
    |\phi ( \delta \bar x + (1-\delta)x) - \delta \phi(\bar x) - (1-\delta) \phi ( x) | &\leq  K \delta (1-\delta) (1+|\bar x|^{p-2} + |x|^{p-2}) |\bar{x}-x|^2, 
    \end{aligned}
    $$
    for any $x, \bar x \in \mathbb R ^n$, $\delta \in [0,1]$. 
    \item\label{codition ellipticity} The matrix $a(x):=\sigma^* \sigma (x) $ is uniformly elliptic; that is, there exists a constant $\kappa_\sigma >0$ such that 
    $$
    \sum_{i,j=1}^n a_{i,j} (x) z_i z_j \geq \kappa_\sigma |z|^2, \quad \text{for any $z \in \mathbb R ^n$.}
    $$
    \item\label{condition rho} $\rho$ is large enough: namely,
    %$\rho > \frac{\hat c _0 ( 4(p-1))}{2}, \hat c _1(2), {\hat c _1(4)} ,     \frac{\hat c _0 ( 2(p-1))+ \hat c _1(2)}{2},     \frac{\hat c _0(2p)}{2},    \hat c _0 (p),    \frac{\hat c _0 ( 2(p-1))+ \hat c _2(2)}{2},    \hat c _2 (2),\\    \frac{\hat c _2 (4)}2,     \frac{\hat c _1 (8)}{2},     \frac{\hat c _0 ( 4(p-2))}{2}    $
     $$
     \begin{aligned}
    \rho > \max \Big\{ & 
    \frac{\hat c _0 ( 4(p-1))}{2}, 
    \hat c _1(2), 
    {\hat c _1(4)} , 
    \frac{\hat c _0 ( 2(p-1))+ \hat c _1(2)}{2}, 
    \frac{\hat c _0(2p)}{2},
     \hat c _0 (p), \\
    & \qquad  \frac{\hat c _0 ( 2(p-1))+ \hat c _2(2)}{2}, 
    \hat c _2 (2),
    \frac{\hat c _2 (4)}2, 
    \frac{\hat c _1 (8)}{2}, 
    \frac{\hat c _0 ( 4(p-2))}{2} \Big\},
    \end{aligned}
    $$
    where, for generic $q \geq 0$ we define the constants
    \begin{equation}\label{eq definition constants}
    \begin{aligned}
        \hat c _0(q) &:= q \Big(\frac32 L + (q-1) L^2 \Big), \\
        \hat c _1(q) &:= q \Big(L+\frac{q-1}{2}L^2 \Big), \\
        \hat c _2 (q)& := {2 q\Big(L+\frac{2 q-1}{2} L^{2}\Big) }+ {(2 q-1) L+2(q-1)^{2} L^{2}}, 
    \end{aligned}
    \end{equation}
    which are related to the growth behavior of the underlying diffusion (see the estimates \eqref{eq estimate SDE growth}-\eqref{eq estimate SDE semiconv SUP} below).
    \end{enumerate}
\end{assumption}

\begin{lemma}
    \label{lemma monotonicity boundary} 
    Under Assumption \ref{assumption general}, for $0\leq \lambda \leq \bar \lambda$, we have that:
    \begin{enumerate}
        \item If $y > e^{-1}$, then  $\mathcal E_{\bar \lambda} (y) \subseteq   \mathcal E_{ \lambda} (y)$;
        \item If $y=e^{-1}$, then $\mathcal E_{\bar \lambda}(y) = \mathcal E_{ \lambda} (y) $;
        \item If $ y<e^{-1}$, then $ \mathcal E_{ \lambda} (y) \subseteq   \mathcal E_{\bar \lambda} (y) $.
    \end{enumerate}
\end{lemma}

\begin{align}
\label{eq estimate SDE growth} 
\mathbb{E}[ |  X ^x_t |^{q} ]  & \leq C (1+|x|^q) e^{  \hat c _0(q) t}, \quad q\geq 1,   \\ 
\label{eq estimate SDE Lip} \mathbb{E}[ |  X^{\bar x}_t - X ^x_t |^{q} ] & \leq |\bar{x} - x|^q e^{  \hat c _1(q) t} , \quad q\geq 2, \\ 
 \label{eq estimate SDE semiconv} 
 \mathbb{E} \big[  \big|  \delta X^{\bar x}_t + (1-\delta) X^x_t - X ^{\delta \bar x + (1-\delta) x}_t \big|^q \big] & \leq C \delta (1-\delta) |\bar{x} - x|^{2q} e^{  \hat c _2(q) t}, \quad q \geq 2,
\end{align}

\begin{align}
\label{eq estimate SDE growth SUP} 
\mathbb{E}\bigg[ \sup_{t \leq \tau}  e^{-\rho t} |  X ^x_t |^{q} \bigg]  & \leq C (1+|x|^q), \quad q = 2(p-1),p,2(p-2),  \\ 
\label{eq estimate SDE Lip SUP} 
\mathbb{E} \bigg[ \sup_{t \leq \tau}  e^{-\rho t} |  X^{\bar x}_t - X ^x_t |^{q} \bigg] & \leq C |\bar{x} - x|^q , \quad q= 2, 4  \\ 
 \label{eq estimate SDE semiconv SUP} 
 \mathbb{E} \bigg[ \sup_{t \leq \tau}  e^{-\rho t} \big|  \delta X^{\bar x}_t + (1-\delta) X^x_t - X ^{\delta \bar x + (1-\delta) x}_t \big|^2 \bigg] & \leq C \delta (1-\delta) |\bar{x} - x|^4,
\end{align}

\begin{lemma}
    \label{lemma limits stopping times}
     Under Assumption \ref{assumption general}, for any $x \in \R ^n$, we have $\tau^\lambda (x,1) \to \tau ^0(x)$, $\P$-a.s., as $\lambda \to 0$.
\end{lemma}

\begin{proof}
    Thanks to Lemma \ref{lemma monotonicity boundary}, 
    for $0\leq \lambda \leq \bar \lambda$, we have that
     $\mathcal E_{\bar \lambda} (1) \subseteq   \mathcal E_{ \lambda} (1)$. 
     This implies that $\mathcal S_{ \lambda} (1) \subseteq   \mathcal S_{\bar \lambda} (1)$, from which we deduce that
     $$
     \tau^{\lambda} (x,1) = \inf \{ t \, | \, X_t^x \in \mathcal S _\lambda (1) \} \geq \inf \{ t \, | \, X_t^x \in \mathcal S _{\bar \lambda} (1) \} = \tau^{\bar \lambda} (x,1).
     $$
 Therefore, the family $\tau^\lambda (x,1)$ is nonincreasing in $\lambda$.
 Using this monotonicity, we can define a limiting stopping time 
 $$
 \tau ^*(x):= \sup_{\lambda >0} \tau^\lambda (x,1) = \lim _{\lambda \to 0} \tau^\lambda (x,1),
 $$
 and we notice that $\tau^0(x) \geq \tau^* (x)$.

Next, by optimality of $\tau^\lambda (x,1)$, we have
$$
\begin{aligned}
    &\mathbb{E}\left[\int_{0}^{\tau^\lambda (x,1)} e ^{-\rho t} \left( \pi\left(X^x_{t}\right) - \lambda (1+\log y) \right) \dd t+e^{-\rho {\tau^\lambda (x,1)}} G\left(X^x_{\tau^\lambda (x,1)}\right)\right] \\
    & \quad \leq  \mathbb{E}\left[\int_{0}^{\tau} e ^{-\rho t} \left( \pi\left(X^x_{t}\right) - \lambda (1+\log y) \right) \dd t+e^{-\rho {\tau}} G\left(X^x_{\tau}\right)\right], \quad \text{for any $\tau \in \mathcal T$.}
\end{aligned}
$$
Taking limits as $\lambda \to 0$ in the latter inequality (using the growth conditions in Assumption \ref{assumption general} and the dominated convergence theorem), we deduce that
$$
\begin{aligned}
    &\mathbb{E}\left[\int_{0}^{\tau^* (x)} e ^{-\rho t}  \pi\left(X^x_{t}\right)  \dd t+e^{-\rho {\tau^* (x)}} G\left(X^x_{\tau^* (x)}\right)\right] \\
    & \quad \leq  \mathbb{E}\left[\int_{0}^{\tau} e ^{-\rho t}   \pi\left(X^x_{t}\right) \dd t+e^{-\rho {\tau}} G\left(X^x_{\tau}\right)\right], \quad \text{for any $\tau \in \mathcal T$.}
\end{aligned}
$$
Hence, $\tau^*(x)$ is also optimal. Thus, since   $\tau^0(x)$ is the minimal optimal stopping time and $\tau^0(x) \geq \tau^* (x)$, we conclude that $\tau^*(x) = \tau^0(x)$. 
\end{proof} 

\end{document}
